\chapter*{ВВЕДЕНИЕ}
\addcontentsline{toc}{chapter}{ВВЕДЕНИЕ}
Основной моделью взаимодействия между двумя и более конечными устройствами в компьютерной сети 
является модель клиент-сервер.
Под термином компьютерная сеть будем понимать набор автономных компьютеров, связанных одной технологией. 
Причём назовём два компьютера связанными, если они могут обмениваться информацией.\cite{tanen_book} 
Таким образом, приложение, реализующее данную модель взаимодействия, 
имеет два типа процессов: 
\begin{itemize}
	\item клиент - это процесс, который обращается к серверу за ресурсом;
	\item сервер - это процесс, который обладает ресуром и обслуживает клиентов.
\end{itemize}
Для обмена информацией приложение должно использовать 
протоколы прикладного уровня, в частности HTTP.
В терминах приложения веб-сервер - это процесс сервер, 
который прослушивает определенный порт для обслуживания клиента,
реализует взаимодействие с клиентом посредством протокола HTTP, и распоряжается, имеющимися ресурсами.
\par
Цель работы - разработка сервера для отдачи статического содержимого с диска по протоколу HTTP.
\par
Для достижения поставленной цели необходимо решить следующие задачи:
\begin{itemize}
	\item проанализировать архитектуру компьютерной сети, протоколы и системные вызовы; 
	\item описать основные компоненты и алгоритмы работы веб-сервера.
	\item реализовать сервер;
	\item провести нагрузочное тестирование разработанного сервера;
\end{itemize}

